
\section{环面与赫尔曼非闭合引理}
\label{sec:torus-nonclosing}

\hzsync{theorem}{10}
\hzsync{equation}{14}

设 $\psi:S\times(-1,+1)\to U\subset M$ 是\trSymplecticManifold{} $(M,\omega)$ 中%
以 $S$ 为模型的一个\trParameterizedHypersurfaceFamily，则 $U$ 是诸\trHypersurface%
$S_{\varepsilon}=\psi\bigl(S\times\{\varepsilon\}\bigr)$ 的并，而且我们知道
\[
c_{0}(U,\omega)<\infty\quad\Longrightarrow\quad \text{存在}\ \varepsilon\ \text{使得}%
\ \mathcal{P}\bigl(S_{\varepsilon}\bigr)\neq\emptyset\ .
\]
反之，当然，若对每个 $\varepsilon$ 都有 $\mathcal{P}\bigl(S_{\varepsilon}\bigr)=\emptyset$，%
则 $c_{0}(U,\omega)=\infty$，从而由\trCapacity{}的\trMonotonicity，%
$c_{0}(M,\omega)=+\infty$。下面我们将描述一个属于 E.~\trZehnder{}[231] 的例子，%
它所阐明的正是这种情形。我们考察\trManifold
\[
M\ =\ T^{3}\times[0,1]\ ,\qquad \text{其中}\ \ T^{3}=\R^{3}/\Z^{3}
\]
即三维\trTorus，并在 $M$ 上通过轻微修改 $\R^{4}$ 上的\trStandard{}结构 $\omega_{0}$%
来定义一个特定的\trSymplecticStructure。回顾一下：对每个满足 $\det A\neq0$ 与%
$A^{T}=-A$ 的常值矩阵 $A$，我们都可以借助
\[
\omega(X,Y)\ =\ \bigl\langle AX,\ Y\bigr\rangle
\]
在 $\R^{4}$ 上配以\trSymplecticStructure{} $\omega$，当然随后也在 $M$ 上如此。%
一个\trHamiltonianFunction{} $H:M\to\R$ 简单地说就是一个函数%
$H:\R^{3}\times[0,1]\to\R$，它关于 $x_{1}$、$x_{2}$、$x_{3}$ 是\trPeriodic{}为 $1$%
的周期函数。\trHamiltonianVectorField{} $X_{H}$ 按通常方式由
\[
\omega\bigl(X_{H},\,Y\bigr)\ =\ -dH(Y)
\]
对一切 $Y\in\R^{4}$ 定义，由此得到
\[
X_{H}\ =\ -A^{-1}\nabla H\ .
\]
现在我们以 $-A^{-1}=J^{*}$ 来定义 $A$，即
\begin{equation}\label{eq:J-star-matrix}
J^{*}\ =\ \begin{pmatrix}
0 & 1 & 0 & \alpha_{1}\\
-1 & 0 & 0 & \alpha_{2}\\
0 & 0 & 0 & 1\\
-\alpha_{1} & -\alpha_{2} & -1 & 0
\end{pmatrix},
\end{equation}
其中 $\alpha_{1}$、$\alpha_{2}\in\R$ 为两个实数。于是 $\det J^{*}=1$，而%
\trSymplectic{}形式 $\omega=\omega^{*}$ 便是下述\trTwoForm：
\begin{equation}\label{eq:omega-star-form}
\omega^{*}\ =\ dx_{2}\wedge dx_{1}\ +\ dx_{4}\wedge dx_{3}\ +\ \alpha_{1}\,dx_{3}\wedge dx_{2}%
\ +\ \alpha_{2}\,dx_{1}\wedge dx_{3}\ .
\end{equation}
选取\trHamiltonianFunction{} $H_{0}:M\to\R$ 为
\begin{equation}\label{eq:H0-coordinate}
H_{0}(x)\ =\ x_{4},
\end{equation}
则 $M$ 上的\trHamiltonianVectorField
\[
X_{H_{0}}\ =\ J^{*}\nabla H_{0}\ =\ J^{*}e_{4}\ =\ \bigl(\alpha_{1},\alpha_{2},1,0\bigr),
\]
是常值的。诸\trEnergySurface
\begin{equation}\label{eq:H0-level-c}
H_{0}^{-1}(c)\ =\ \Bigl\{\,x\in M\ \Bigm|\ H_{0}(x)\ =\ x_{4}\ =\ c\,\Bigr\}
\end{equation}
都是\trRegular{}的，并且由于 $x_{j}\in S^{1}$（$1\leq j\leq 3$），它们恰是三维%
\trTorus{} $H_{0}^{-1}(c)=T^{3}\times\{c\}$。我们引入缩记
\begin{equation}\label{eq:alpha-vector}
\alpha\ =\ \bigl(\alpha_{1},\alpha_{2},1\bigr)\in\R^{3}\ .
\end{equation}
限制在诸\trEnergySurface{} $T^{3}\times\{c\}$ 上的\trHamiltonianEquation{}之\trFlow{}是%
\trLinear{}的，并且由
\begin{equation}\label{eq:linear-flow}
\dot{x}_{k}\ =\ \alpha_{k}\ ,\qquad 1\leq k\leq 3\ .
\end{equation}
确定。若我们选取\trRationallyIndependent{}的 $\alpha$，即要求
\[
\langle\alpha,j\rangle\ \neq\ 0\qquad\text{对}\quad 0\neq j\in\Z^{3},
\]
则依据\trKronecker{}定理，\eqref{eq:linear-flow} 的一切\trOrbit{}都在 $T^{3}$ 上%
稠密。于是，我们便得到一个\trHamiltonianSystem{}的例子：它的诸\trEnergySurface{}都是%
\trRegular{}且\trCompact{}的，然而它却不容许任何\trPeriodicSolution。顺便指出，%
这些\trEnergySurface{}都不是\trContactType{}的。

顺带地，我们指出：人们可以容易地构造一个\trEmbedding{}%
$\varphi:T^{3}\times[0,1]\to\R^{4}$。以 $T^{3}$ 为模型的、由它诱导出的那个%
\trParameterizedHypersurfaceFamily，相对于被\trPushforward{}的\trSymplecticStructure%
（我们仍把它记作 $\omega^{*}$）不带任何\trClosedCharacteristic。然而，%
$\varphi(M)\subset\R^{4}$ 上的 $\omega^{*}$ 并不等价于从 $\R^{4}$ 诱导出的%
\trStandard{}结构 $\bigl(\varphi(M),\omega_{0}\bigr)$。的确，$\omega_{0}$ 是一个%
\trExact{}的\trSymplecticStructure，而 $\omega^{*}$ 却不是；此外，%
$c_{0}\bigl(\varphi(M),\omega^{*}\bigr)=+\infty$ 而%
$c_{0}\bigl(\varphi(M),\omega_{0}\bigr)<\infty$。特别地，这个例子并不能充当对下述%
猜想的一个\trCounterexample：$(\R^{2n},\omega_{0})$ 中每张\trHypersurface{}都具有%
一条\trClosedCharacteristic。我们指出，由 \eqref{eq:linear-flow} 在 $T^{3}$ 上%
定义的\trVectorField{} $X$ 满足：对 $T^{3}$ 上\trClosed{}的 1-形式 $\beta=dx_{1}$%
有 $\beta(X)\neq0$。现在更一般地假设：$S$ 是任意一个无\trBoundary{}的\trCompact%
\trManifold{}，$\dim S=2n-1$，而 $X$ 是 $S$ 上的一个\trVectorField，对 $S$ 上某个%
\trClosed{}的 1-形式 $\beta$ 满足 $\beta(X)\neq0$。那么，$S$ 不能作为一张%
\trHypersurface{} $S\subset(M,\omega)$ 嵌入任何一个具有\trExact{}\trSymplecticStructure%
（即 $\omega=d\lambda$）的\trSymplecticManifold{}中，使得 $X\in\mathcal{L}_{S}$，%
即 $X$ 属于\trCharacteristicLineBundle。的确，倘若它能够如此，那么由于%
$\beta(X)\neq0$ 且 $X\in\ker d\lambda|_{S}$，我们便可断定%
$\beta\wedge(d\lambda)^{n-1}$ 是 $S$ 上的一个\trVolumeForm，从而
\[
\operatorname{vol}(S)\ =\ \int_{S}\beta\wedge(d\lambda)^{n-1}\ \neq\ 0\ .
\]
然而这是不可能的，因为形式 $\beta\wedge(d\lambda)^{n-1}$ 等于%
$d\bigl(\lambda\wedge\beta\wedge(d\lambda)^{n-2}\bigr)$，从而是 $S$ 上一个%
\trExact{}形式。

现在我们对\trSymplecticStructure{} $\omega^{*}$ 作进一步的限制：要求 $\alpha$ 不仅%
\trRationallyIndependent，而且还满足\trDiophantine{}条件
\begin{align*}
\text{(D.C.)}\qquad & \ \bigl\lvert\langle\alpha,j\rangle\bigr\rvert\ \geq\ \gamma\,\lvert j\rvert^{-\tau}
\end{align*}
对一切 $0\neq j\in\Z^{3}$ 成立，这里 $\gamma>0$ 与 $\tau>2$ 为两个常数。%
\trAlmostEvery{}向量 $\alpha\in\R^{3}$ 都满足这样的条件。事实表明，上述例子在%
$H_{0}$ 的、具有充分多阶导数的小扰动之下实际上是\trDynamics{}稳定的。这一令人惊讶的%
现象是 M.~\trHerman{}[110, 111] 最近发现的，并参看 J.C.~\trYoccoz{}[229]。在%
我们的框架下他证明了

\begin{theorem}[M.~\trHerman{}的\trNonClosingLemma]\label{thm:herman-nonclosing}
考察\trSymplecticManifold{} $\bigl(M,\omega^{*}\bigr)=\bigl(T^{3}\times[0,1],\omega^{*}\bigr)$，%
其\trSymplecticStructure{}为 \eqref{eq:omega-star-form}，而 $\alpha\in\R^{3}$ 满足%
条件 (D.C.)。令 $H_{0}(x)=x_{4}$。那么，存在 $H_{0}$ 的一个\trNeighborhood{} $W$
\[
W\ =\ \Bigl\{H\in C^{k}(M)\ \Bigm|\ \bigl\lvert H-H_{0}\bigr\rvert_{C^{k}}\ \leq\ \delta\Bigr\},
\]
（这里 $\delta>0$ 充分小而 $k>\tau+2$），使得对每个 $H\in W$ 与每个满足%
$\tfrac{1}{2}\leq c\leq\tfrac{3}{2}$ 的 $c$，\trRegular\trEnergySurface
\[
H^{-1}(c)\ \subset\ M
\]
都不具有 $X_{H}$ 的\trPeriodicSolution。
\end{theorem}

我们指出，这种\trDynamics\trRigidity{}现象是那个特殊的\trSymplecticStructure{} $\omega^{*}$%
的性质。我们把证明推后，先指出\trHerman{}定理的一个重要推论。

依据\trPoincare{}\trRecurrence{}定理（见第 1 章），对一个\trHamiltonianSystem{} $H$%
而言，\trCompact\trEnergySurface{}上的\trAlmostEvery{}点 $x$ 都是\trRecurrence{}点，%
于是人们会问：是否存在着一个附近的系统，它有一条经过 $x$ 附近的%
\trPeriodic\trTrajectory。这就是所谓的\trClosingProblem。答案取决于对何为「附近」的%
理解。在其工作 [179]（1983 年）中，C.C.~\trPugh{}与 C.~\trRobinson{}对%
\trHamiltonian{}情形证明了：存在一列\trHamiltonianSystem{} $H_{j}$ 以及这些近似%
系统的一列相应的\trPeriodic{}点 $x_{j}$，满足
\[
H_{j}\ \to\ H\quad\text{在}\ C^{2}\ \text{中}\qquad\text{以及}\qquad x_{j}\ \to\ x\ .
\]
这是文献 [184] 中那个著名的 $C^{1}$-\trClosingLemma{}的一个特殊情形。至于能否选取%
诸 $H_{j}$ 使得对 $k>2$ 有 $H_{j}\to H$ 于 $C^{k}$ 中，这是一个长期悬而未决的%
公开问题。M.~\trHerman{}对这个问题给出了令人惊讶的回答：在特殊\trManifold{}%
$(M,\omega^{*})$ 上，只要 $k$ 充分大——大致地说，大于所论\trManifold{}的维数——%
$C^{k}$-\trClosingLemma{}便不成立。

\begin{proof}[\cref*{thm:herman-nonclosing}的证明]
考察 $M$ 上由
\[
H\ =\ x_{4}+h(x)\qquad\text{以及}\qquad \lvert h\rvert_{C^{k}}\ \ \text{很小}
\]
给出的 $H_{0}$ 的那个\trNeighborhood。\trHamiltonianVectorField{} $X_{H}$ 可表示为
\begin{equation}\label{eq:XH-components}
X_{H}(x)\ =\ J^{*}\nabla H(x)\ =\
\begin{pmatrix}
H_{x_{2}}\ +\ \alpha_{1}H_{x_{4}}\\
-H_{x_{1}}\ +\ \alpha_{2}H_{x_{4}}\\
H_{x_{4}}\\
-H_{x_{3}}\ -\ \alpha_{1}H_{x_{1}}\ -\ \alpha_{2}H_{x_{2}}
\end{pmatrix}.
\end{equation}
由
\[
x_{4}+h(x)\ =\ H(x)\ =\ c,
\]
所定义的关于 $c$ 的\trEnergySurface{}是一张\trEmbedding{}的\trTorus{} $T^{3}$。%
的确，我们可以对 $x_{4}$ 求解，并得到 $x_{4}=K\bigl(x_{1},x_{2},x_{3},c\bigr)$，%
其中函数 $K$ 满足恒等式
\begin{equation}\label{eq:K-identity}
H\Bigl(x_{1},\,x_{2},\,x_{3},\,K\bigl(x_{1},x_{2},x_{3},c\bigr)\Bigr)\ =\ c
\end{equation}
对一切 $x\in T^{3}$ 成立。这些函数关于 $x\in T^{3}$ 是周期的。为了计算%
$H^{-1}(c)$ 上的\trVectorField{} $X_{H}$，我们微分恒等式 \eqref{eq:K-identity}，%
得到 $H_{x_{j}}=-H_{x_{4}}\cdot K_{x_{j}}$ 在 $H^{-1}(c)$ 上对 $1\leq j\leq 3$%
成立。于是，由 \eqref{eq:XH-components}，\trVectorField{} $X_{H}$ 在 $H^{-1}(c)$%
上便由
\begin{equation}\label{eq:XH-rho}
\begin{pmatrix}\dot{x}_{1}\\[2pt] \dot{x}_{2}\\[2pt] \dot{x}_{3}\end{pmatrix}
\ =\ \rho\cdot
\begin{pmatrix}\alpha_{1}\ -\ K_{x_{2}}\\[2pt] \alpha_{2}\ +\ K_{x_{1}}\\[2pt] 1\end{pmatrix}
\qquad\text{在 $T^{3}$ 上},
\end{equation}
给出，其中 $\rho(x)=H_{x_{4}}(x,K)$ 是一个正函数。\trVectorField{} \eqref{eq:XH-rho}%
经\trReparametrization{}后与
\begin{equation}\label{eq:V-field}
\begin{pmatrix}\dot{x}_{1}\\[2pt] \dot{x}_{2}\\[2pt] \dot{x}_{3}\end{pmatrix}
\ =\ V(x)\ :=\
\begin{pmatrix}\alpha_{1}\ -\ K_{x_{2}}\\[2pt] \alpha_{2}\ +\ K_{x_{1}}\\[2pt] 1\end{pmatrix}
\qquad\text{在 $T^{3}$ 上}\ .
\end{equation}
具有相同的\trOrbit。注意，$T^{3}$ 上的这个\trVectorField{} $V$ 是 $T^{3}$ 上常值%
\trVectorField{} $\dot{x}=\alpha$ 的一个\trPerturbation，而且我们可以应用 K.A.M.%
理论的一个特殊情形。我们将利用 J.~\trMoser{}于 1966 年在 [161] 中给出的下述结果。

\begin{theorem}\label{thm:moser-small-denominator}
在\trTorus{} $T^{n}=\R^{n}/\Z^{n}$ 上考察形如
\[
\dot{x}\ =\ \alpha+f(x),
\]
的一个\trVectorField，并假设常值向量 $\alpha\in\R^{n}$ 满足\trDiophantine{}条件
\[
\bigl\lvert\langle\alpha,j\rangle\bigr\rvert\ \geq\ \gamma\lvert j\rvert^{-\tau},
\]
对一切 $0\neq j\in\Z^{n}$，其中 $\gamma>0$ 与 $\tau>n-1$ 为两个常数。那么，若%
$\lvert f\rvert_{C^{m}}$ 充分小且 $m>\tau+1$，则存在常值向量 $\lambda\in\R^{n}$%
与一个在\trIdentityMap{}附近的 $C^{1}$-\trDiffeomorphism{} $u=id+v:T^{n}\to T^{n}$，%
使得
\[
(du)^{-1}\bigl(\alpha+f+\lambda\bigr)\circ u\ =\ \alpha\ .
\]
\end{theorem}

换言之，如果一个\trVectorField{}连同充分多阶导数都靠近常值\trVectorField{} $\alpha$，%
那么它可以通过常数 $\lambda\in\R^{n}$ 加以修改，使得这个修改后的\trVectorField{}%
等价于 $\alpha$。这个开创性的\trSmallDenominator{}结果，其证明基于一种在无穷多个%
空间中进行的、分析上十分精微的迭代技巧，此处不再展开；我们请读者参看 [161]。

一般地说 $\lambda\neq0$，这一点由 $f$ 取常值的例子即已表明。但在我们所论的情形中，%
\trVectorField{} $f$ 具有由问题的\trHamiltonian{}性质遗传下来的一个附加结构，正是%
这个附加结构使我们得以断定 $\lambda=0$。这将由下述引理得到。

\textbf{\trhLemma{}。}\quad 设 $V$ 是\trTorus{} $T^{n}$ 上一个满足%
$\operatorname{div}V=0$ 的\trVectorField。假设
\[
(du)^{-1}\cdot V\circ u\ =\ \alpha,
\]
其中 $\alpha\in\R^{n}$ 是一个有理无关向量，而 $u=id+v$ 是\trTorus{}的一个%
\trDiffeomorphism。那么，$u$ 保持 $T^{n}$ 上的\trLebesgue{}测度，并且
\[
\int_{T^{n}}V\ =\ \alpha\ .
\]

我们把引理的证明推后，先证明\trHerman{}的结果。把 \eqref{eq:V-field} 右端的%
\trVectorField{}缩记为 $\alpha+f(x)$，我们首先由\cref{thm:moser-small-denominator}%
得出：存在 $\lambda\in\R^{3}$，使得\trVectorField
\[
V\ :=\ \alpha+f(x)+\lambda
\]
被变换为 $T^{3}$ 上的常值\trVectorField{} $\dot{x}=\alpha$。\trVectorField{} $f$%
具有由\trHamiltonian{}结构遗传下来的一种特殊形式：即 $\int f=0$，因为周期函数的%
导数的积分恒为零。从而
\[
\int_{T^{3}}V\ =\ \alpha+\lambda\ .
\]
此外，$V$ 还满足 $\operatorname{div}V=0$。于是，由该引理我们得出%
$\alpha+\lambda=\alpha$，从而 $\lambda=0$。我们已经证明：\trVectorField{}%
\eqref{eq:V-field} 在 $T^{3}$ 上的\trFlow{}共轭于 $\dot{x}=\alpha$ 的%
\trLinear\trFlow，因而不具有任何\trPeriodicSolution。这对每个满足%
$\tfrac{1}{2}\leq c\leq\tfrac{3}{2}$ 的常数 $c$ 都成立\footnote{原书证明结尾印作「$\tfrac{1}{2}\leq c\leq\tfrac{2}{3}$」，与其定理陈述（原书 p.139）之「$\tfrac{1}{2}\leq c\leq\tfrac{3}{2}$」相抵，此处依定理陈述统一为 $\tfrac{3}{2}$。——译者注}。%
\cref{thm:herman-nonclosing}的证明至此完成，剩下的只是证明引理。

积分 $V\circ u=du(\alpha)$ 并注意到 $u=id+v$，我们得到
\[
\int_{T^{n}}V\circ u\ d\mu\ =\ \alpha
\]
因为 $dv(\alpha)$ 的积分消失。为了证明引理，我们只须证明 $u^{*}\mu=\mu$，即 $u$%
保持\trLebesgue{}测度。\trVectorField{} $V$ 的\trFlow{} $\psi^{t}$ 满足
\[
\bigl(\psi^{t}\bigr)^{*}\mu\ =\ \mu\ .
\]
这是假设 $\operatorname{div}V=0$ 的一个推论。此外，
\[
u^{-1}\circ\psi^{t}\circ u\ =\ R_{t\alpha}\ ,\qquad t\in\R
\]
其中 $R_{t\alpha}:x\mapsto x+t\alpha$ 是\trTorus{}上的\trTranslation。从而%
$\psi^{t}\circ u=u\circ R_{t\alpha}$，于是由 $\bigl(\psi^{t}\bigr)^{*}\mu=\mu$ 我们%
得到
\[
u^{*}\mu\ =\ \bigl(R_{t\alpha}\bigr)^{*}\bigl(u^{*}\mu\bigr)\ ,\qquad t\in\R\ .
\]
由假设，$\alpha$ 是有理无关的，故一切\trOrbit{} $t\mapsto R_{t\alpha}(x)$ 在 $T^{n}$%
上稠密。由于 $\mu$ 是一个常值形式，我们得出 $u^{*}\mu=a\cdot\mu$，其中 $a$ 为%
常数。在\trTorus{}上积分
\[
\int_{T^{n}}\mu\ =\ \int_{T^{n}}u^{*}\mu\ =\ a\int_{T^{n}}\mu
\]
便得 $a=1$，这就证明了确实有 $u^{*}\mu=\mu$。引理至此证毕。
\end{proof}

显然，\trSymplecticManifold{} $\bigl(T^{4},\omega^{*}\bigr)$ 也具有无穷的%
\trCapacity{} $c_{0}$
\[
c_{0}\bigl(T^{4},\omega^{*}\bigr)\ =\ +\infty
\]
只要 $\omega^{*}$ 是由 \eqref{eq:omega-star-form} 所定义的那个特殊的%
\trSymplecticStructure，且 $\alpha=\bigl(\alpha_{1},\alpha_{2},1\bigr)$ 取为有理无关的。%
我们知道，在二维情形，对每个\trVolumeForm{} $\omega$ 都有%
$c_{0}\bigl(T^{2},\omega\bigr)<\infty$。至于当 $n>1$ 时
\[
c_{0}\bigl(T^{2n},\omega_{0}\bigr)\ <\ \infty
\]
是否成立，这是一个公开问题。这里 $\omega_{0}$ 是从 $\R^{2n}$ 诱导出的%
\trStandard\trSymplecticStructure。
